\section{问题重述}

\subsection{问题背景}

干燥是决定中药材成品品质的关键工序之一，热风烘干是常见方式。
该方式主要包括\textbf{预热平衡}和\textbf{恒温干燥}两个阶段，
通过调控烘房温湿环境完成药材干燥。
工艺参数选取不当会导致干燥效率低、能耗高、成品品质不稳定。
传统试验优化模式成本高、周期长，需要借助数理分析与数值仿真得到干燥规律。

\subsection{问题提出}

\begin{enumerate}[leftmargin=2.2em, itemsep=2pt]
  \item \textbf{问题 1}：建立\textbf{预热平衡阶段}药材温度和水分浓度变化规律的
        数学模型，给出指定时刻、指定位置的结果，并保存完整结果。
  \item \textbf{问题 2}：建立\textbf{整个烘干过程}药材温度和水分浓度变化规律的
        数学模型，给出指定时刻、指定位置的结果，并保存完整结果。
  \item \textbf{问题 3}：按烘干要求（药材\textbf{各处}水分浓度低于
        $0.15$ kg/kg），确定药材烘干所需时间。
  \item \textbf{问题 4}：考虑药材因水分流失发生\textbf{尺寸变化}，
        确定药材的烘干时长。
\end{enumerate}

\subsection{已知条件}

药材为圆柱体，长 $L_0=0.25$ m、初始半径 $R_0=0.02$ m；
初始温度 $28\ ^\circ\mathrm{C}$、初始干基含水率 $2.55$ kg/kg。
对流换热系数 $h=25$ W/(m$^2\cdot$K)、对流传质系数 $h_m=8\times10^{-7}$ m/s。
题目给出三套物性经验公式：
附录2（问题1，常物性）、附录3（问题2/3，变物性）、附录4（问题4，变物性）；
附件1 给出烘房温度与水分浓度的实测序列（241 点，0--14400 s），
附件2 给出药材半径随时间的变化（145 点，0--259200 s）。
